\begin{mistake}{2026-05-07}{01}

\begin{problemcontent}
已知平面上三条不同直线的方程分别为
\[
l_1:ax+2by+3c=0,\qquad
l_2:bx+2cy+3a=0,\qquad
l_3:cx+2ay+3b=0.
\]
试证：这三条直线交于一点的充分必要条件为
\[
a+b+c=0.
\]
\end{problemcontent}

\begin{solutioncontent}
先证充分性。若 \(a+b+c=0\)，取点
\[
P\left(3,\frac32\right).
\]
代入三条直线，分别有
\[
3a+3b+3c=3(a+b+c)=0,
\]
\[
3b+3c+3a=3(a+b+c)=0,
\]
\[
3c+3a+3b=3(a+b+c)=0.
\]
故三条直线均经过 \(P\left(3,\frac32\right)\)，充分性成立。

再证必要性。设三条直线交于公共点 \((x_0,y_0)\)，则
\[
\begin{pmatrix}
a&2b&3c\\
b&2c&3a\\
c&2a&3b
\end{pmatrix}
\begin{pmatrix}
x_0\\y_0\\1
\end{pmatrix}
=
\begin{pmatrix}
0\\0\\0
\end{pmatrix}.
\]
由于 \((x_0,y_0,1)^T\neq 0\)，齐次线性方程组有非零解，因此系数行列式为零：
\[
\det
\begin{pmatrix}
a&2b&3c\\
b&2c&3a\\
c&2a&3b
\end{pmatrix}=0.
\]
计算该行列式：
\[
\begin{aligned}
D
&=a
\begin{vmatrix}
2c&3a\\
2a&3b
\end{vmatrix}
-2b
\begin{vmatrix}
b&3a\\
c&3b
\end{vmatrix}
+3c
\begin{vmatrix}
b&2c\\
c&2a
\end{vmatrix}\\
&=a(6bc-6a^2)-2b(3b^2-3ac)+3c(2ab-2c^2)\\
&=18abc-6(a^3+b^3+c^3)\\
&=-6(a^3+b^3+c^3-3abc)\\
&=-6(a+b+c)(a^2+b^2+c^2-ab-bc-ca).
\end{aligned}
\]
于是
\[
(a+b+c)(a^2+b^2+c^2-ab-bc-ca)=0.
\]
又
\[
a^2+b^2+c^2-ab-bc-ca
=\frac12\left[(a-b)^2+(b-c)^2+(c-a)^2\right].
\]
若第二个因子为零，则 \(a=b=c\)。此时三条方程相同；若 \(a=b=c\neq 0\)，三条直线重合，不是三条不同直线；若 \(a=b=c=0\)，方程退化为 \(0=0\)，也不是直线。均与题设矛盾。

故只能有
\[
a+b+c=0.
\]
综上，三条不同直线交于一点的充分必要条件为
\[
\boxed{a+b+c=0}.
\]
\end{solutioncontent}

\mistakeknowledge{
\begin{itemize}
  \item \textbf{核心公式/定理}：三条直线 \(A_i x+B_i y+C_i=0\) 若交于有限点，则 \(\det(A_i,B_i,C_i)_{3\times 3}=0\)。本题反向还需结合“三条不同直线”排除退化情形。
  \item \textbf{易错点}：把三式相加只能得到 \((a+b+c)(x_0+2y_0+3)=0\)，不能直接推出 \(a+b+c=0\)。
  \item \textbf{关联知识}：恒等式 \(a^3+b^3+c^3-3abc=(a+b+c)(a^2+b^2+c^2-ab-bc-ca)\)，且 \(a^2+b^2+c^2-ab-bc-ca=\frac12[(a-b)^2+(b-c)^2+(c-a)^2]\)。
\end{itemize}}

\end{mistake}
